% \iffalse meta-comment
%
% hit-spread.dtx 是字间距撑开
%
% \fi
%
% \section{字间距撑开}
%
% \changes{v3.2a}{2026/09/27}{定宽拉开在 \hologo{pdfLaTeX} 下按字拆（一个汉字是几个记号，逐记号拆会把字劈开）}
% \changes{v3.2a}{2026/09/26}{\option{heading-spread} 的默认值照范例改成两个半角空格，按平衡
%   选项排成半格；取值以控制空格打头时 \cs{hit@spread@by:nn} 插它本身}
% \changes{v3.2a}{2026/08/28}{加 \cs{hit@spread@by:nn} 与 \cs{hit@spread@to:nn}：标题的
%   字间距从常量取值里抽出来（撑开写进取值会漏进目录条目与 PDF 书签），封面字段标签
%   分散对齐到固定宽度，底层统一走 \env{CJKfilltwosides}；两个参数都 \texttt{f} 展开
%   （桶到 \texttt{base} 的回落是一层转指，只展开一次拿到的还是个宏）}
%
% 在排版时添加字间距，词条本身保留原文以供书签使用。
% 参数须用 f 展开解析多层引用；单次展开无法取得内容，完全展开又会执行排版命令。
% 间距为空时原样输出；控制空格按半角空格规则处理，长度值使用 hspace。
%
%    \begin{macrocode}
%<*shared-spread>
%    \end{macrocode}
%
%    \begin{macrocode}
\tl_new:N \l__hit_spread_tl
\tl_new:N \l__hit_spread_gap_tl
\bool_new:N \l__hit_spread_first_bool
\tl_const:Ne \c__hit_spread_cspace_tl { \exp_not:c { ~ } }
\cs_new_eq:Nc \__hit_spread_cspace: { ~ }
\cs_new_protected:Npn \hit@spread@by:nn #1#2
  {
    \tl_set:Nf \l__hit_spread_gap_tl {#1}
    \tl_if_blank:VTF \l__hit_spread_gap_tl {#2}
      {
        \tl_if_head_eq_meaning:VNTF \l__hit_spread_gap_tl \__hit_spread_cspace:
          {
            \__hit_spread_sbcs:
            \__hit_spread_by:Vn \l__hit_spread_gap_tl {#2}
          }
          { \__hit_spread_by:nn { \hspace {#1} } {#2} }
      }
  }
\cs_new_protected:Npn \__hit_spread_sbcs:
  {
    \cs_if_exist:NT \docxlike_format_sbcs_space:
      {
        \bool_if:NT \g_docxlike_format_balance_bool
          {
            \exp_args:NNV \tl_replace_all:Nnn \l__hit_spread_gap_tl
              \c__hit_spread_cspace_tl { \docxlike_format_sbcs_space: }
          }
      }
  }
\cs_new_protected:Npn \__hit_spread_by:nn #1#2
  {
    \tl_set:Nf \l__hit_spread_tl {#2}
    \bool_set_true:N \l__hit_spread_first_bool
    \tl_map_inline:Nn \l__hit_spread_tl
      {
        \bool_if:NTF \l__hit_spread_first_bool
          { \bool_set_false:N \l__hit_spread_first_bool }
          {#1}
        ##1
      }
  }
\cs_generate_variant:Nn \__hit_spread_by:nn { V }
\seq_new:N \l__hit_spread_seq
\sys_if_engine_pdftex:T
  {
    \cs_set_protected:Npn \__hit_spread_by:nn #1#2
      {
        \hit@engine@pdftex@chars:nN {#2} \l__hit_spread_seq
        \seq_use:Nn \l__hit_spread_seq {#1}
      }
  }
%    \end{macrocode}
%
% 按指定宽度分散对齐，只伸展汉字间胶，不在标点边界额外插入间距。
%    \begin{macrocode}
\cs_new_protected:Npn \hit@spread@to:nn #1#2
  { \hit@engine@spread@to:nn {#1} {#2} }
%    \end{macrocode}
%
% \env{CJKfilltwosides} 是个 \env{minipage}，内容比给的宽度宽就折行。要的是「压进
% 这个宽度、不折行」的话它做不到，那种情况仍写 \cs{makebox}\oarg{宽度}\oarg{s}。
% \file{src/config/hit-terms.dtx} 里 \option{major-field-label} 的
% \texttt{学科/专业学位类别} 是唯一一处：九个字压进八个字宽。
%
% 配置层的取值是在 \pkg{expl3} 之外记号化的（见 CONTRIBUTING §4.7），那里 \texttt{:}
% 不是字母，\cs{hit@spread@to:nn} 会被拆成 \cs{hit@spread@to} 加 \texttt{:nn}，
% 报 Undefined control sequence。所以两个宏各留一个不带冒号的名字，取值里写这个：
% \begin{latex}
% supervisor-field-label = {\hit@spread@to{6em}{导师}}
% \end{latex}
% 与 \cs{hit@if@option@TF} 同一个路数——凡是要在排版内容里出现的，名字里都不带冒号。
%    \begin{macrocode}
\cs_new_eq:NN \hit@spread@to \hit@spread@to:nn
\cs_new_eq:NN \hit@spread@by \hit@spread@by:nn
%    \end{macrocode}
%
% \cs{hit@sp} 是原件里在汉字中间敲的一个半角空格，给不能逐字撑开的标签用（研究生报告
% 封面的“学\ \ \ \ \ 院（部）”只在头一个字后面敲了五个）。规矩与上面的控制空格一样。
%    \begin{macrocode}
\cs_new_protected:Npn \hit@sp
  {
    \bool_lazy_and:nnTF
      { \cs_if_exist_p:N \docxlike_format_sbcs_space: }
      { \g_docxlike_format_balance_bool }
      { \docxlike_format_sbcs_space: }
      { \c_space_tl }
  }
%    \end{macrocode}
%
%    \begin{macrocode}
%</shared-spread>
%    \end{macrocode}
%
% \Finale
